\documentclass[11pt]{article}

\usepackage[T1]{fontenc}
\usepackage[utf8]{inputenc}
\usepackage{amsmath,amssymb,amsthm}
\usepackage{geometry}
\usepackage{hyperref}
\geometry{margin=1in}

\hypersetup{
  colorlinks=true,
  linkcolor=blue,
  urlcolor=blue,
  citecolor=blue
}

\newtheorem{theorem}{Theorem}
\newtheorem{lemma}[theorem]{Lemma}
\newtheorem{corollary}[theorem]{Corollary}
\newtheorem{proposition}[theorem]{Proposition}
\theoremstyle{definition}
\newtheorem{definition}[theorem]{Definition}
\theoremstyle{remark}
\newtheorem{remark}[theorem]{Remark}

\newcommand{\Z}{\mathbb Z}
\newcommand{\N}{\mathbb N}
\newcommand{\Zi}{\mathbb Z[i]}
\newcommand{\Norm}{\operatorname N}
\newcommand{\vp}{v_p}

\title{A Complete Integer Solution of\\[3pt]
\texorpdfstring{$y^2+x^3y+z^2+1=0$}{y2+x3y+z2+1=0}}
\author{}
\date{}

\begin{document}
\maketitle

\begin{abstract}
We solve the Diophantine equation
\[
        y^2+x^3y+z^2+1=0
\]
completely over the integers.  The solution is reduced, by completing the square, to representations of
$x^6-4$ as a sum of two squares.  For completeness, the classical two-squares criterion is proved from first principles using the Euclidean arithmetic of the Gaussian integers.
\end{abstract}

\section{Statement of the answer}

For a positive integer $n$, define
\[
\mathcal R_n
=
\{(a,b)\in \Z_{>0}^2: a^2+b^2=n^6-4,\ a\equiv 1\pmod 2,\ b\equiv0\pmod2\}.
\]
Also define
\[
\mathcal N
=
\{n\in\Z_{>0}: n\equiv3\pmod {12}\text{ and }\vp(n^6-4)\text{ is even for every prime }p\equiv3\pmod4\}.
\]

\begin{theorem}[Complete solution]
The integer solutions of
\[
        y^2+x^3y+z^2+1=0
\]
are exactly the triples
\[
\boxed{
        x=\varepsilon n,
        \qquad
        y=\frac{\sigma a-\varepsilon n^3}{2},
        \qquad
        z=\tau\frac b2
}
\]
where
\[
        n\in\mathcal N,
        \qquad
        (a,b)\in\mathcal R_n,
        \qquad
        \varepsilon,\sigma,\tau\in\{\pm1\}.
\]
Equivalently, for a positive integer $n$, one has $\mathcal R_n\ne\varnothing$ if and only if
\[
        n\equiv3\pmod {12}
\]
and every prime $p\equiv3\pmod4$ occurs to an even exponent in $n^6-4$.
\end{theorem}

The proof occupies the rest of the note.  The formulation above is finite and effective: for a given $n$, factor $n^6-4$ and check the exponents of the primes congruent to $3$ modulo $4$.

\section{Elementary reduction to a sum of two squares}

\begin{lemma}[The square-completion identity]
For integers $x,y,z$, the equation
\[
        y^2+x^3y+z^2+1=0
\]
is equivalent to
\[
        (2y+x^3)^2+(2z)^2=x^6-4.
\]
\end{lemma}

\begin{proof}
Multiplying the original equation by $4$ gives
\[
        4y^2+4x^3y+4z^2+4=0.
\]
Since
\[
        4y^2+4x^3y=(2y+x^3)^2-x^6,
\]
this is equivalent to
\[
        (2y+x^3)^2-x^6+4z^2+4=0,
\]
which is precisely
\[
        (2y+x^3)^2+(2z)^2=x^6-4.
\]
The algebra is reversible, so the two equations are equivalent.
\end{proof}

\begin{lemma}[First congruence obstructions]
If $(x,y,z)\in\Z^3$ satisfies
\[
        y^2+x^3y+z^2+1=0,
\]
then $x$ is odd and $3\mid x$.
\end{lemma}

\begin{proof}
First suppose that $x$ is even.  Then $x^3y$ is divisible by $4$, and the original equation gives
\[
        y^2+z^2+1\equiv0\pmod4.
\]
But a square modulo $4$ is either $0$ or $1$, so $y^2+z^2$ is congruent only to $0,1$, or $2$ modulo $4$, never to $3$ modulo $4$.  This contradiction proves that $x$ is odd.

Now suppose that $3\nmid x$.  By the square-completion identity,
\[
        (2y+x^3)^2+(2z)^2=x^6-4.
\]
The quadratic residues modulo $9$ are
\[
        0,1,4,7.
\]
Indeed, this is obtained by squaring the residue classes $0,1,\ldots,8$.  No sum of two elements of $\{0,1,4,7\}$ is congruent to $6$ modulo $9$: the possible sums are
\[
        0,1,2,4,5,7,8 \pmod9.
\]
On the other hand, if $3\nmid x$, then $x^6\equiv1\pmod9$, and hence
\[
        x^6-4\equiv1-4\equiv6\pmod9.
\]
Thus the right-hand side of the square-completion identity is congruent to $6$ modulo $9$, while the left-hand side is a sum of two squares modulo $9$, impossible.  Therefore $3\mid x$.
\end{proof}

\begin{lemma}[Parity of representations]
Let $n$ be odd and suppose
\[
        a^2+b^2=n^6-4
\]
with $a,b\in\Z$.  Then neither $a$ nor $b$ is zero.  Moreover, exactly one of $a,b$ is odd and the other is even.  After interchanging the two coordinates if necessary, the representation may be written with $a$ odd and $b$ even.
\end{lemma}

\begin{proof}
Since $n$ is odd, $n^2\equiv1\pmod8$, so $n^6\equiv1\pmod8$.  Hence
\[
        n^6-4\equiv5\pmod8.
\]
No square is congruent to $5$ modulo $8$, so neither coordinate can be zero.  Also, because $n^6-4$ is odd, exactly one of $a,b$ is odd and the other is even.  Therefore we can label the odd coordinate by $a$ and the even coordinate by $b$.
\end{proof}

\section{The two-squares theorem, proved self-containedly}

We need the classical criterion for an integer to be a sum of two squares.  To keep the solution self-contained, we prove it here.

\subsection{Gaussian integers}

\begin{definition}
The Gaussian integers are
\[
        \Zi=\{u+vi:u,v\in\Z\}.
\]
For $\alpha=u+vi\in\Zi$, define its norm by
\[
        \Norm(\alpha)=\alpha\overline{\alpha}=u^2+v^2.
\]
\end{definition}

\begin{lemma}[Multiplicativity of the norm]
For all $\alpha,\beta\in\Zi$,
\[
        \Norm(\alpha\beta)=\Norm(\alpha)\Norm(\beta).
\]
\end{lemma}

\begin{proof}
Using complex conjugation,
\[
        \Norm(\alpha\beta)
        =\alpha\beta\overline{\alpha\beta}
        =\alpha\beta\overline\alpha\overline\beta
        =(\alpha\overline\alpha)(\beta\overline\beta)
        =\Norm(\alpha)\Norm(\beta).
\]
\end{proof}

\begin{lemma}[Euclidean division in $\Zi$]
For any $\alpha,\beta\in\Zi$ with $\beta\ne0$, there exist $q,r\in\Zi$ such that
\[
        \alpha=q\beta+r
        \qquad\text{and}\qquad
        \Norm(r)<\Norm(\beta).
\]
\end{lemma}

\begin{proof}
Write $\alpha/\beta=s+ti$ with $s,t\in\mathbb R$.  Choose integers $m,n$ satisfying
\[
        |s-m|\le\frac12,
        \qquad
        |t-n|\le\frac12,
\]
and set $q=m+ni\in\Zi$.  Define $r=\alpha-q\beta$.  Then
\[
        r=\beta\left(\frac{\alpha}{\beta}-q\right),
\]
and therefore
\[
\Norm(r)
=\Norm(\beta)\left|\frac{\alpha}{\beta}-q\right|^2
\le
\Norm(\beta)\left(\frac14+\frac14\right)
=\frac12\Norm(\beta)<\Norm(\beta).
\]
Thus the desired division exists.
\end{proof}

\begin{lemma}[$\Zi$ is a unique factorization domain]
Every nonzero nonunit Gaussian integer factors as a product of irreducible Gaussian integers, and this factorization is unique up to multiplication by units and reordering.
\end{lemma}

\begin{proof}
The preceding Euclidean division gives the Euclidean algorithm in $\Zi$.  Hence for any $\alpha,\beta\in\Zi$, not both zero, there is a greatest common divisor $d$ and Gaussian integers $u,v$ such that
\[
        d=u\alpha+v\beta.
\]
This is obtained exactly as over $\Z$: repeatedly divide with remainder of strictly smaller norm until the last nonzero remainder appears, and then back-substitute.

We next show that every irreducible element of $\Zi$ is prime.  Let $\pi\in\Zi$ be irreducible and suppose
\[
        \pi\mid \alpha\beta.
\]
Let $d$ be a greatest common divisor of $\pi$ and $\alpha$.  Since $d\mid\pi$ and $\pi$ is irreducible, either $d$ is a unit or $d$ is associated to $\pi$.  If $d$ is associated to $\pi$, then $\pi\mid\alpha$.  If $d$ is a unit, Bezout's identity gives
\[
        u\pi+v\alpha=1
\]
for some $u,v\in\Zi$.  Multiplying by $\beta$ gives
\[
        u\pi\beta+v\alpha\beta=\beta.
\]
Both terms on the left are divisible by $\pi$, because $\pi\mid\alpha\beta$.  Hence $\pi\mid\beta$.  Thus an irreducible element dividing a product divides one of the factors; that is, every irreducible is prime.

Existence of factorizations follows by descent on the positive integer norm.  If a nonzero nonunit $\gamma\in\Zi$ has no factorization into irreducibles, choose such a counterexample with $\Norm(\gamma)$ minimal.  Then $\gamma$ is not irreducible, so
\[
        \gamma=\delta\eta
\]
with $\delta,\eta$ nonunits.  Hence
\[
        1<\Norm(\delta)<\Norm(\gamma),
        \qquad
        1<\Norm(\eta)<\Norm(\gamma).
\]
By minimality, both $\delta$ and $\eta$ factor into irreducibles, and therefore so does $\gamma$, a contradiction.

Uniqueness is now standard and follows from primality of irreducibles.  If
\[
        \pi_1\cdots\pi_r=\rho_1\cdots\rho_s
\]
are two factorizations into irreducibles, then $\pi_1$ divides the product on the right, so $\pi_1$ divides some $\rho_j$.  Since $\rho_j$ is irreducible, $\pi_1$ and $\rho_j$ are associates.  Cancel them and repeat by induction.  This proves uniqueness up to units and reordering.
\end{proof}

\subsection{\texorpdfstring{Rational primes in $\Zi$}{Rational primes in Z[i]}}

\begin{lemma}[Fermat's little theorem]
Let $p$ be a rational prime and let $u\in\Z$ with $p\nmid u$.  Then
\[
        u^{p-1}\equiv1\pmod p.
\]
\end{lemma}

\begin{proof}
Multiplication by $u$ permutes the nonzero residue classes modulo $p$.  Therefore
\[
        u\cdot 2u\cdot 3u\cdots (p-1)u
        \equiv
        1\cdot2\cdot3\cdots(p-1)
        \pmod p.
\]
The right-hand product is not divisible by $p$, so it may be cancelled modulo $p$, giving
\[
        u^{p-1}\equiv1\pmod p.
\]
\end{proof}

\begin{lemma}[Wilson's theorem]
If $p$ is a rational prime, then
\[
        (p-1)!\equiv-1\pmod p.
\]
\end{lemma}

\begin{proof}
Modulo $p$, each nonzero residue has a multiplicative inverse.  The only residues equal to their own inverse are the solutions of
\[
        t^2\equiv1\pmod p,
\]
namely $t\equiv1$ and $t\equiv-1\pmod p$.  Therefore, in the product of all nonzero residues modulo $p$, every residue except $1$ and $-1$ pairs with its distinct inverse, and each such pair contributes $1$.  Thus
\[
        (p-1)!\equiv 1\cdot(-1)\equiv-1\pmod p.
\]
\end{proof}

\begin{lemma}[The residue of $-1$ modulo an odd prime]
Let $p$ be an odd rational prime.
\begin{enumerate}
\item If $p\equiv3\pmod4$, then $-1$ is not a quadratic residue modulo $p$.
\item If $p\equiv1\pmod4$, then $-1$ is a quadratic residue modulo $p$.
\end{enumerate}
\end{lemma}

\begin{proof}
First let $p\equiv3\pmod4$.  If $t^2\equiv-1\pmod p$, then $p\nmid t$, so Fermat's little theorem gives $t^{p-1}\equiv1\pmod p$.  But $(p-1)/2$ is odd, and therefore
\[
        t^{p-1}=(t^2)^{(p-1)/2}\equiv(-1)^{(p-1)/2}\equiv-1\pmod p,
\]
a contradiction.  Hence $-1$ is not a quadratic residue modulo $p$.

Now let $p\equiv1\pmod4$.  Wilson's theorem gives
\[
        (p-1)!\equiv-1\pmod p.
\]
Pairing the factors $j$ and $p-j$ for $1\le j\le (p-1)/2$ gives
\[
        (p-1)!
        =\prod_{j=1}^{(p-1)/2}j(p-j)
        \equiv
        (-1)^{(p-1)/2}\left(\left(\frac{p-1}{2}\right)!\right)^2
        \pmod p.
\]
Since $p\equiv1\pmod4$, the exponent $(p-1)/2$ is even.  Therefore
\[
        \left(\left(\frac{p-1}{2}\right)!\right)^2\equiv-1\pmod p,
\]
so $-1$ is a quadratic residue modulo $p$.
\end{proof}

\begin{lemma}[Behavior of rational primes in $\Zi$]
Let $p$ be a rational prime.
\begin{enumerate}
\item $2=(1+i)(1-i)$, so $2$ is a norm: $2=\Norm(1+i)$.
\item If $p\equiv3\pmod4$, then $p$ is irreducible, hence prime, in $\Zi$.
\item If $p\equiv1\pmod4$, then there are integers $a,b$ such that
\[
        p=a^2+b^2=\Norm(a+bi).
\]
Equivalently, $p$ factors nontrivially in $\Zi$.
\end{enumerate}
\end{lemma}

\begin{proof}
The identity for $2$ is immediate:
\[
        (1+i)(1-i)=2.
\]

Let $p\equiv3\pmod4$.  Suppose $p$ were reducible in $\Zi$.  Then
\[
        p=\alpha\beta
\]
with $\alpha,\beta$ nonunits.  Taking norms gives
\[
        p^2=\Norm(\alpha)\Norm(\beta).
\]
Because $\alpha$ and $\beta$ are nonunits, their norms exceed $1$.  Since $p$ is rationally prime, this forces
\[
        \Norm(\alpha)=\Norm(\beta)=p.
\]
Writing $\alpha=a+bi$, we obtain
\[
        a^2+b^2=p.
\]
If $b=0$, then $a^2=p$, impossible.  Thus $b\not\equiv0\pmod p$, and reducing modulo $p$ gives
\[
        \left(ab^{-1}\right)^2\equiv -1\pmod p,
\]
contradicting the fact that $-1$ is not a quadratic residue modulo primes $p\equiv3\pmod4$.  Therefore $p$ is irreducible in $\Zi$.  Since $\Zi$ is a unique factorization domain, irreducible elements are prime, so $p$ is prime in $\Zi$.

Now let $p\equiv1\pmod4$.  By the preceding lemma, choose $t\in\Z$ such that
\[
        t^2\equiv-1\pmod p.
\]
Then
\[
        p\mid t^2+1=(t+i)(t-i)
\]
in $\Zi$.  If $p$ were prime in $\Zi$, it would divide $t+i$ or $t-i$.  But $p\mid t+i$ in $\Zi$ would mean that both real and imaginary parts of $t+i$ are divisible by $p$, i.e. $p\mid t$ and $p\mid1$, impossible.  Similarly, $p\nmid t-i$.  Hence $p$ is not prime in $\Zi$, and therefore not irreducible.  So
\[
        p=\alpha\beta
\]
with $\alpha,\beta$ nonunits.  Taking norms again gives
\[
        p^2=\Norm(\alpha)\Norm(\beta),
\]
and hence $\Norm(\alpha)=p$ and $\Norm(\beta)=p$.  Writing $\alpha=a+bi$ gives
\[
        p=a^2+b^2.
\]
This proves the claim.
\end{proof}

\subsection{The criterion}

\begin{theorem}[Two-squares criterion]
Let $M$ be a positive integer.  Then $M$ is representable as
\[
        M=A^2+B^2
\]
with $A,B\in\Z$ if and only if every rational prime $p\equiv3\pmod4$ occurs to an even exponent in the prime factorization of $M$.
\end{theorem}

\begin{proof}
First suppose $M=A^2+B^2$.  Put
\[
        \alpha=A+Bi\in\Zi.
\]
Then
\[
        M=\Norm(\alpha)=\alpha\overline{\alpha}.
\]
Let $p\equiv3\pmod4$ be a rational prime.  By the preceding lemma, $p$ is prime in $\Zi$.  In the unique factorization of $\alpha$, suppose $p$ occurs to exponent $e$.  Since $\overline p=p$, the same exponent $e$ occurs in the unique factorization of $\overline\alpha$.  Therefore the exponent of $p$ in
\[
        \alpha\overline\alpha=M
\]
is $2e$, an even integer.  Hence every rational prime $p\equiv3\pmod4$ has even exponent in $M$.

Conversely, suppose every prime $p\equiv3\pmod4$ has even exponent in $M$.  Write the rational prime factorization of $M$ as
\[
        M=2^e
        \prod_{p\equiv1\, (4)}p^{e_p}
        \prod_{q\equiv3\, (4)}q^{2f_q}.
\]
For $2$, use $2=\Norm(1+i)$.  For each prime $p\equiv1\pmod4$, the preceding lemma gives a Gaussian integer $\pi_p$ with
\[
        \Norm(\pi_p)=p.
\]
For each prime $q\equiv3\pmod4$, the Gaussian integer $q$ itself has norm
\[
        \Norm(q)=q^2.
\]
Define
\[
        \gamma=(1+i)^e
        \prod_{p\equiv1\, (4)}\pi_p^{e_p}
        \prod_{q\equiv3\, (4)}q^{f_q}
        \in\Zi.
\]
By multiplicativity of the norm,
\[
        \Norm(\gamma)=M.
\]
Writing $\gamma=A+Bi$ gives
\[
        M=A^2+B^2.
\]
This proves both directions.
\end{proof}

\section{Completion of the Diophantine solution}

We now return to
\[
        y^2+x^3y+z^2+1=0.
\]

\begin{lemma}[Necessary form of $|x|$]
Let $(x,y,z)$ be an integer solution, and put $n=|x|$.  Then
\[
        n\equiv3\pmod {12},
\]
and every prime $p\equiv3\pmod4$ occurs to an even exponent in $n^6-4$.
\end{lemma}

\begin{proof}
By the congruence lemma, $x$ is odd and $3\mid x$.  Hence $n$ is odd and $3\mid n$.  By the square-completion identity,
\[
        (2y+x^3)^2+(2z)^2=n^6-4.
\]
Thus $n^6-4$ is a sum of two squares.  By the two-squares criterion, every prime $p\equiv3\pmod4$ occurs to an even exponent in $n^6-4$.

It remains to prove $n\equiv3\pmod4$.  Since $3\mid n$ and $n>0$, we have $n\ge3$.  Suppose, for contradiction, that $n\equiv1\pmod4$.  Then
\[
        n^3-2\equiv3\pmod4,
        \qquad
        n^3+2\equiv3\pmod4.
\]
Moreover,
\[
        \gcd(n^3-2,n^3+2)\mid4.
\]
Both $n^3-2$ and $n^3+2$ are odd, so in fact
\[
        \gcd(n^3-2,n^3+2)=1.
\]
Now an odd integer congruent to $3$ modulo $4$ must contain some prime congruent to $3$ modulo $4$ to an odd exponent: indeed, primes congruent to $1$ modulo $4$ contribute $1$ modulo $4$, while primes congruent to $3$ modulo $4$ contribute $-1$ modulo $4$, so the total number of prime factors congruent to $3$ modulo $4$, counted with multiplicity, must be odd.

Therefore each of the coprime integers $n^3-2$ and $n^3+2$ contains a prime $\equiv3\pmod4$ to an odd exponent.  Since the two factors are coprime, that prime remains to an odd exponent in their product
\[
        (n^3-2)(n^3+2)=n^6-4,
\]
contradicting the two-squares criterion.  Hence $n\not\equiv1\pmod4$.  Since $n$ is odd, $n\equiv3\pmod4$.

Combining $3\mid n$ with $n\equiv3\pmod4$ gives
\[
        n\equiv3\pmod {12}.
\]
This proves the lemma.
\end{proof}

\begin{lemma}[Sufficiency and reconstruction]
Let $n\in\mathcal N$, let $(a,b)\in\mathcal R_n$, and choose signs
\[
        \varepsilon,\sigma,\tau\in\{\pm1\}.
\]
Define
\[
        x=\varepsilon n,
        \qquad
        y=\frac{\sigma a-\varepsilon n^3}{2},
        \qquad
        z=\tau\frac b2.
\]
Then $x,y,z$ are integers and
\[
        y^2+x^3y+z^2+1=0.
\]
Conversely, every integer solution arises in this way.
\end{lemma}

\begin{proof}
First observe that $x$ is an integer.  Since $n\equiv3\pmod {12}$, $n$ is odd, so $n^3$ is odd.  Since $a$ is odd, $\sigma a-\varepsilon n^3$ is even; hence $y\in\Z$.  Since $b$ is even, $z\in\Z$.

Now compute
\[
        2y+x^3
        =\sigma a-\varepsilon n^3+(\varepsilon n)^3.
\]
Because $\varepsilon^3=\varepsilon$, this simplifies to
\[
        2y+x^3=\sigma a.
\]
Similarly,
\[
        2z=\tau b.
\]
Therefore
\[
        (2y+x^3)^2+(2z)^2=a^2+b^2=n^6-4=x^6-4.
\]
By the square-completion identity, this is equivalent to
\[
        y^2+x^3y+z^2+1=0.
\]
Thus every triple constructed in the stated way is a solution.

Conversely, let $(x,y,z)$ be any integer solution and set
\[
        n=|x|,
        \qquad
        \varepsilon=\operatorname{sgn}(x).
\]
The necessary-form lemma gives $n\in\mathcal N$.  Define
\[
        A=2y+x^3,
        \qquad
        B=2z.
\]
Then
\[
        A^2+B^2=n^6-4.
\]
Since $x$ is odd, $A=2y+x^3$ is odd, and $B=2z$ is even.  The parity lemma shows that neither $A$ nor $B$ is zero.  Put
\[
        a=|A|,
        \qquad
        b=|B|,
        \qquad
        \sigma=\operatorname{sgn}(A),
        \qquad
        \tau=\operatorname{sgn}(B).
\]
Then $(a,b)\in\mathcal R_n$, and
\[
        2y+x^3=A=\sigma a.
\]
Since $x=\varepsilon n$, this gives
\[
        2y=\sigma a-\varepsilon n^3,
        \qquad
        y=\frac{\sigma a-\varepsilon n^3}{2}.
\]
Also
\[
        z=\tau\frac b2.
\]
Thus the solution has exactly the asserted form.
\end{proof}

\begin{lemma}[Non-emptiness of $\mathcal R_n$]
For a positive integer $n$, the set $\mathcal R_n$ is nonempty if and only if
\[
        n\equiv3\pmod {12}
\]
and every prime $p\equiv3\pmod4$ occurs to an even exponent in $n^6-4$.
\end{lemma}

\begin{proof}
Suppose first that $\mathcal R_n$ is nonempty, and choose $(a,b)\in\mathcal R_n$.  Then $a^2+b^2=n^6-4$, so the two-squares criterion gives the required even-exponent condition.  Also, reducing $a^2+b^2=n^6-4$ modulo $2$ gives $1\equiv n^6\pmod2$, so $n$ is odd.  Therefore
\[
        x=n,
        \qquad
        y=\frac{a-n^3}{2},
        \qquad
        z=\frac b2
\]
are integers, and the square-completion identity gives
\[
        y^2+x^3y+z^2+1=0.
\]
Applying the necessary-form lemma to this solution gives $n\equiv3\pmod {12}$.

Conversely, suppose $n\equiv3\pmod {12}$ and every prime $p\equiv3\pmod4$ occurs to an even exponent in $n^6-4$.  By the two-squares criterion, there exist integers $c,d$ such that
\[
        c^2+d^2=n^6-4.
\]
Because $n$ is odd,
\[
        n^6-4\equiv5\pmod8.
\]
The parity lemma then shows that neither $c$ nor $d$ is zero and exactly one of them is odd.  Taking absolute values and, if necessary, interchanging the two coordinates, we obtain positive integers $a,b$ with
\[
        a^2+b^2=n^6-4,
        \qquad
        a\text{ odd},
        \qquad
        b\text{ even}.
\]
Thus $(a,b)\in\mathcal R_n$, so $\mathcal R_n\ne\varnothing$.
\end{proof}

The complete theorem follows from the reconstruction lemma and the non-emptiness lemma.

\section{Examples}

For $n=3$, one has
\[
        3^6-4=729-4=725.
\]
The positive representations with the first coordinate odd and the second even are
\[
        725=25^2+10^2=23^2+14^2=7^2+26^2.
\]
Taking $x=3$ gives the following solutions:
\[
\begin{array}{c|c|c}
(a,b)& y & z \\
\hline
(25,10)& -1,-26 & \pm5 \\
(23,14)& -2,-25 & \pm7 \\
(7,26)& -10,-17 & \pm13
\end{array}
\]
Thus, for example,
\[
        (x,y,z)=(3,-1,5)
\]
is a solution, since
\[
        (-1)^2+3^3(-1)+5^2+1=1-27+25+1=0.
\]
The symmetry
\[
        (x,y,z)\mapsto(-x,-y,z)
\]
preserves the equation, so the corresponding solutions with $x=-3$ are obtained by changing the sign of $y$.

A larger example is obtained from
\[
        87^6-4=653251^2+83002^2.
\]
Taking
\[
        x=87,
        \qquad
        a=653251,
        \qquad
        b=83002,
        \qquad
        \sigma=1,
        \qquad
        \tau=1
\]
gives
\[
        y=\frac{653251-87^3}{2}
        =\frac{653251-658503}{2}
        =-2626,
        \qquad
        z=41501.
\]
Hence
\[
        (87,-2626,41501)
\]
is another integer solution.

\section{Concluding form}

The equation
\[
        y^2+x^3y+z^2+1=0
\]
has no sporadic solutions outside the family described above.  Every solution is governed by a representation
\[
        |x|^6-4=a^2+b^2
\]
with $a$ odd and $b$ even, and the exact admissibility condition on $|x|$ is the two-squares condition
\[
        \vp(|x|^6-4)\equiv0\pmod2
        \qquad\text{for every prime }p\equiv3\pmod4,
\]
together with
\[
        |x|\equiv3\pmod {12}.
\]
This gives a full unconditional classification of the integer solutions.

\end{document}
